% !TeX program = pdflatex
% Cennik wynagrodzeń za rozwiązania projektowe --- Basilisk Systems P.S.A. (wzór uchwały Rady Dyrektorów).
% Do Umowy nr 2 § 6 ust. 1 (basilisk_ip.tex).
\documentclass[12pt,a4paper]{report}
\usepackage[T1]{fontenc}
\usepackage[utf8]{inputenc}
\usepackage[english,polish]{babel}
\usepackage{basilisk_i18n}
\usepackage{lmodern}
\usepackage{microtype}
\usepackage{setspace}
\usepackage{parskip}
\usepackage{needspace}
\usepackage{tabularx}

\BusinessPlanCompany{Basilisk Systems}
\BusinessPlanWatermarkText{Basilisk Systems}
\BusinessPlanWatermarkColor{black!6}
\BusinessPlanWatermarkFontSize{38}{46}
\BusinessPlanWatermarkAngle{38}
\BusinessPlanAuthorsText{Basilisk Systems --- Rada Dyrektorów}
\BusinessPlanConfidentialText{Wzór uchwały --- poufne}
\BusinessPlanFooterText{Basilisk Systems --- cennik wynagrodzeń za rozwiązania projektowe (wzór)}
\BusinessPlanVersionText{Wersja 1.0 --- projekt}

\setstretch{1.18}
\setlist[enumerate]{topsep=0.25em,itemsep=0.18em,parsep=0.1em,leftmargin=2.2em}
\setlength{\parindent}{0pt}
\setlength{\parskip}{0.55em}
\setlength{\emergencystretch}{4em}
\renewcommand{\arraystretch}{1.25}

\newcounter{paragraf}
\newcommand{\paragraf}[1]{%
  \refstepcounter{paragraf}%
  \phantomsection%
  \addcontentsline{toc}{section}{\S~\theparagraf\ #1}%
  \Needspace{6\baselineskip}%
  \subsection*{\S~\theparagraf\ #1}%
}
\newenvironment{ContractClause}[1]{\paragraf{#1}}{\par}

\newcommand{\field}[1]{\textbf{[#1]}}
\newcommand{\CompanyFull}{Basilisk Systems prosta spółka akcyjna}

\hypersetup{
  pdftitle={Basilisk Systems P.S.A. --- cennik wynagrodzeń za rozwiązania projektowe},
  pdfauthor={Basilisk Systems --- Rada Dyrektorów}
}

\begin{document}

\BusinessPlanTitlePage
  {Basilisk Systems}
  {Wersja 1.0 --- projekt}
  {Cennik}
  {Cennik wynagrodzeń za rozwiązania projektowe}
  {Wynalazki, wzory użytkowe i wzory przemysłowe twórców współpracujących ze Spółką}
  {1.0}

\tableofcontents
\newpage

\chapter*{Uchwała nr \field{numer} Rady Dyrektorów}
\addcontentsline{toc}{chapter}{Uchwała Rady Dyrektorów --- Cennik}

\begin{center}
Rady Dyrektorów \CompanyFull{} z siedzibą w Pruszkowie\\
z dnia \field{data}\\
w sprawie ustalenia cennika wynagrodzeń za rozwiązania projektowe
\end{center}

Rada Dyrektorów, działając na podstawie §~6 ust.~1 umowy przeniesienia praw do utworów tworzonych dla Spółki (Umowa nr~2) oraz art.~22 ust.~1 ustawy z dnia 30 czerwca 2000 r. --- Prawo własności przemysłowej, ustala, co następuje:

\begin{ContractClause}{Zakres}
\begin{enumerate}
  \item Cennik określa wynagrodzenie twórców wynalazków, wzorów użytkowych i wzorów przemysłowych, do których prawo do uzyskania patentu, prawa ochronnego albo prawa z rejestracji przysługuje Spółce na podstawie Umowy nr~2 (\textbf{„Twórca”}, \textbf{„Rozwiązanie”}).
  \item Cennika nie stosuje się do Twórców będących Założycielami albo dyrektorami Spółki; ich wynagrodzenie jest zawarte w wynagrodzeniu z tytułu stosunku podstawowego i w udziale w Spółce.
  \item Cennik może być załączony do Umowy nr~2 albo doręczony Twórcy przed jej zawarciem; doręczenie może nastąpić w formie dokumentowej.
\end{enumerate}
\end{ContractClause}

\begin{ContractClause}{Stawki}
\begin{enumerate}
  \item Za każde Rozwiązanie Spółka wypłaca jednorazowo:

\begin{center}
\begin{tabularx}{\linewidth}{@{}>{\raggedright\arraybackslash}X >{\raggedleft\arraybackslash}p{3.6cm} >{\raggedleft\arraybackslash}p{3.6cm}@{}}
\toprule
\textbf{Rodzaj Rozwiązania} & \textbf{po dokonaniu zgłoszenia} & \textbf{po uzyskaniu prawa} \\
\midrule
wynalazek (patent) & \field{kwota} zł & \field{kwota} zł \\
wzór użytkowy (prawo ochronne) & \field{kwota} zł & \field{kwota} zł \\
wzór przemysłowy (prawo z rejestracji) & \field{kwota} zł & \field{kwota} zł \\
\bottomrule
\end{tabularx}
\end{center}

  \item Kwoty są kwotami brutto; Spółka pobiera z nich należne zaliczki na podatek i składki.
  \item Wynagrodzenie po dokonaniu zgłoszenia należy się raz za Rozwiązanie, bez względu na liczbę zgłoszeń w innych państwach albo w trybie regionalnym lub międzynarodowym. Wynagrodzenie po uzyskaniu prawa należy się raz, za pierwsze prawo udzielone na to Rozwiązanie.
  \item Jeżeli Spółka korzysta z Rozwiązania bez dokonania zgłoszenia, w szczególności utrzymując je jako tajemnicę przedsiębiorstwa, wynagrodzenie po dokonaniu zgłoszenia wypłaca się w terminie 30 dni od decyzji Rady Dyrektorów o odstąpieniu od zgłoszenia.
\end{enumerate}
\end{ContractClause}

\begin{ContractClause}{Współtwórcy}
Jeżeli Rozwiązanie ma kilku Twórców, wynagrodzenie dzieli się między nich w częściach wskazanych w zgłoszeniu Rozwiązania Spółce, a w braku takiego wskazania --- w częściach równych. Udział twórcy, do którego Cennik nie ma zastosowania (§~1 ust.~2), nie zwiększa udziałów pozostałych Twórców.
\end{ContractClause}

\begin{ContractClause}{Wypłata}
Wynagrodzenie wypłaca się przelewem na rachunek wskazany przez Twórcę w terminie 30 dni od dnia dokonania zgłoszenia w Urzędzie Patentowym Rzeczypospolitej Polskiej albo innym właściwym urzędzie, albo od dnia doręczenia Spółce decyzji o udzieleniu prawa. Prawo do wynagrodzenia nie wygasa z chwilą zakończenia współpracy Twórcy ze Spółką.
\end{ContractClause}

\begin{ContractClause}{Waloryzacja}
Kwoty określone w §~2 ust.~1 podlegają waloryzacji z dniem 1 marca każdego roku, począwszy od roku następującego po roku podjęcia niniejszej uchwały, o średnioroczny wskaźnik cen towarów i usług konsumpcyjnych ogółem za poprzedni rok kalendarzowy, ogłaszany przez Prezesa Głównego Urzędu Statystycznego w Dzienniku Urzędowym Rzeczypospolitej Polskiej „Monitor Polski”, jeżeli wskaźnik ten przekracza 100. Zwaloryzowane kwoty zaokrągla się w górę do pełnych 10 zł. Wynagrodzenie ustala się według kwot obowiązujących w dniu zdarzenia, od którego jest należne. Waloryzacja nie wymaga zmiany Cennika; Spółka udostępnia Twórcom aktualne kwoty.
\end{ContractClause}

\begin{ContractClause}{Zmiany i wejście w życie}
\begin{enumerate}
  \item Rada Dyrektorów może zmienić Cennik uchwałą. Zmianę doręcza się Twórcom w formie dokumentowej.
  \item Zmiana Cennika na niekorzyść Twórcy nie dotyczy Rozwiązań zgłoszonych Spółce przed doręczeniem zmiany.
  \item Uchwała wchodzi w życie z dniem \field{data}.
\end{enumerate}
\end{ContractClause}

\vspace{1.5em}
\field{podpisy dyrektorów}

\end{document}
